% !TeX root = ../se200_identity_regimes.tex

\section{Regime Structure and Neutral Representation}
\label{sec:representation}

This section consolidates the regime profiles established by the preceding
derivations and states the current lower bound on the number of distinct
referential regimes required.

%--------------------------------------------
\subsection{Previously Established Regimes}

The following regime profiles were defined in
Section~\ref{sec:preliminaries} and are not refined by the transformations
analyzed in Sections~\ref{subsec:enr_split}--\ref{subsec:nor_split}:

\begin{itemize}
\item $\mathrm{OBL}$: obligation-bearing entities,
\item $\mathrm{OCC}$: time-indexed occurrences,
\item $\mathrm{REC}$: descriptive records.
\end{itemize}


These profiles do not exhibit split pressure under the transformation
families considered thus far. Their classifications remain internally
consistent with respect to identity preservation.


%--------------------------------------------
\subsection{Refined Regimes}

The preceding sections establish the following refinements:

\begin{itemize}
\item $\mathrm{ENR} \;\Rightarrow\; \mathrm{ENR\mbox{-}L}, \mathrm{ENR\mbox{-}I}$,
\item $\mathrm{CTX} \;\Rightarrow\; \mathrm{CTX\mbox{-}E}, \mathrm{CTX\mbox{-}S}$,
\item $\mathrm{NOR} \;\Rightarrow\; \mathrm{NOR\mbox{-}C}, \mathrm{NOR\mbox{-}S}$.
\end{itemize}

Each refinement is forced by split pressure under a specific transformation:

\begin{itemize}
\item $\mathrm{ENR}$ splits under branching ($\mathrm{BF}$),
\item $\mathrm{CTX}$ splits under decomposition ($\mathrm{AD}$),
\item $\mathrm{NOR}$ splits under refinement ($\mathrm{RF}$).
\end{itemize}

%--------------------------------------------
\subsection{Derived Regime Set}
\label{subsec:derived_regime_set}

This subsection fixes the regime profiles derived in
Sections~\ref{subsec:enr_split}, \ref{subsec:ctx_split}, and \ref{subsec:nor_split}.
These profiles form the basis for the lower bound and pairwise non-collapse results.

\begin{definition}[Derived regime set]
\label{def:derived_regime_set}
The derived regime set consists of the following nine regime profiles:
\[
\{
\mathrm{OBL},\,
\mathrm{OCC},\,
\mathrm{REC},\;
\mathrm{ENR\mbox{-}L},\,
\mathrm{ENR\mbox{-}I},\,
\mathrm{CTX\mbox{-}E},\,
\mathrm{CTX\mbox{-}S},\,
\mathrm{NOR\mbox{-}C},\;
\mathrm{NOR\mbox{-}S}
\}.
\]
\end{definition}

Each regime profile is determined by its identity carrier
and its classification behavior over the fixed transformation basis.

\begin{itemize}
  \item $\mathrm{OBL}$: obligation-bearing entity profile.
  \item $\mathrm{OCC}$: time-indexed occurrence profile.
  \item $\mathrm{REC}$: descriptive record profile.
  \item $\mathrm{ENR\mbox{-}L}$: enduring referent (locus-bound).
  \item $\mathrm{ENR\mbox{-}I}$: enduring referent (instrument-bound).
  \item $\mathrm{CTX\mbox{-}E}$: applicability context (extension-sensitive).
  \item $\mathrm{CTX\mbox{-}S}$: applicability context (structure-sensitive).
  \item $\mathrm{NOR\mbox{-}C}$: normative structure (composition-sensitive).
  \item $\mathrm{NOR\mbox{-}S}$: normative structure (substitution-sensitive).
\end{itemize}


These profiles are defined entirely by their identity and persistence
behavior under the transformation basis.
No additional ontological categories are introduced.

\begin{convention}
  $\mathcal{S}$ denotes an ontology known to be admissible.
  $\mathcal{O}$ denotes an arbitrary ontology whose admissibility
  is under consideration in a given proof or argument.
\end{convention}


%--------------------------------------------
\subsection{Regime-Typed Multigraph Embedding}
\label{subsec:refinement_construction}

\begin{definition}[Regime-typed directed multigraph]
\label{def:regime_multigraph}
Let $\mathcal{S}$ be an admissible ontology with derived regime set
$\mathcal{P}$. The regime-typed directed multigraph
$G(\mathcal{S}) = (V, E, \lambda, \rho)$ is defined by:
\begin{itemize}
  \item $V = \bigsqcup_{P \in \mathcal{P}} \mathcal{E}_P$,
        the disjoint union of equivalence classes over all regime profiles,
  \item $E \subseteq V \times \mathcal{A} \times V$,
        where $\mathcal{A}$ is the admissible relation vocabulary,
  \item $\lambda : V \to \mathcal{P}$ assigns each vertex its regime type,
  \item $\rho : E \to \mathcal{A}$ labels each edge with its admissible relation.
\end{itemize}
\end{definition}

\begin{theorem}[Faithful embedding]
\label{thm:embedding}
The map $\iota : \mathcal{S} \to G(\mathcal{S})$ sending each
identity carrier $x$ to its equivalence class $[x]_P$
under its unique regime profile $P$ (Lemma~\ref{lem:regime_determinacy})
is injective up to regime-typed isomorphism and preserves admissible structure.
\end{theorem}

\begin{proof}
Injectivity follows from regime determinacy
(Lemma~\ref{lem:regime_determinacy}): distinct identity behaviors
map to distinct equivalence classes.
Admissible structure is preserved because edges in $G(\mathcal{S})$
are drawn exclusively from the admissible relation vocabulary $\mathcal{A}$,
and no causal or normative relations are introduced.
Regime-typed isomorphism is the appropriate equivalence
because identity is defined at the level of equivalence classes,
not individual representations.
\end{proof}

This embedding is the construction asserted in the abstract and
consequences sections. Path families over $G(\mathcal{S})$ constitute
the representational basis for structural comparison and transformation.
Full development of path grammar and expressive adequacy
is given in Section~\ref{sec:formal_setting}.


%--------------------------------------------
\subsection{Lower Bound}

\begin{theorem}[Current lower bound]
Any substrate satisfying the requirements of
Sections~\ref{sec:preliminaries}--\ref{subsec:nor_split}
must realize at least nine pairwise non-collapsing regime profiles.
\end{theorem}

\begin{proof}
Sections~\ref{subsec:enr_split}, \ref{subsec:ctx_split}, and \ref{subsec:nor_split}
establish non-collapse for the pairs
$(\mathrm{ENR\mbox{-}L}, \mathrm{ENR\mbox{-}I})$,
$(\mathrm{CTX\mbox{-}E}, \mathrm{CTX\mbox{-}S})$, and
$(\mathrm{NOR\mbox{-}C}, \mathrm{NOR\mbox{-}S})$.

These refinements establish nine distinct profiles:
$\mathrm{ENR\mbox{-}L}$, $\mathrm{ENR\mbox{-}I}$,
$\mathrm{CTX\mbox{-}E}$, $\mathrm{CTX\mbox{-}S}$,
$\mathrm{NOR\mbox{-}C}$, and $\mathrm{NOR\mbox{-}S}$.

Together with the previously established profiles
$\mathrm{OBL}$, $\mathrm{OCC}$, and $\mathrm{REC}$,
this yields nine candidate regime profiles.

Pairwise non-collapse for this regime set is established in
Section~\ref{subsec:pairwise_noncollapse}.
Therefore the substrate must realize at least nine pairwise non-collapsing regime profiles.
\end{proof}

%--------------------------------------------
\subsection{Pairwise Non-Collapse Matrices}
\label{subsec:pairwise_noncollapse}

To organize the pairwise non-collapse argument, we record the status of each unordered pair
of derived regime profiles in matrix form.
The first matrix gives the coarse status of each pair:
entries marked $\mathrm{THM}$ are already separated by prior split results,
while entries marked $\mathrm{GEN}$ are separated by a generic cross-family criterion.
The next two matrices refine the generic entries by indicating the kind of separation used.

\paragraph{Generic separation criteria.}
For distinct profiles $P,Q$, a matrix entry marked GEN indicates that
$P$ and $Q$ differ in at least one of the following regime-defining respects:
(i) carrier of identity,
(ii) identity basis,
(iii) applicability classification over the fixed transformation basis,
or (iv) persistence behavior.
Any such difference that changes the induced equivalence relation yields non-collapse.

\medskip
\noindent\textbf{Coarse pairwise non-collapse matrix.}

\[
\begin{array}{@{}c|ccccccccc@{}}
 & \mathrm{OBL} & \mathrm{OCC} & \mathrm{REC} & \mathrm{ENR\mbox{-}L} & \mathrm{ENR\mbox{-}I} & \mathrm{CTX\mbox{-}E} & \mathrm{CTX\mbox{-}S} & \mathrm{NOR\mbox{-}C} & \mathrm{NOR\mbox{-}S} \\
\hline
\mathrm{OBL} & -   & GEN & GEN & GEN & GEN & GEN & GEN & GEN & GEN \\
\mathrm{OCC} & GEN & -   & GEN & GEN & GEN & GEN & GEN & GEN & GEN \\
\mathrm{REC} & GEN & GEN & -   & GEN & GEN & GEN & GEN & GEN & GEN \\
\mathrm{ENR\mbox{-}L} & GEN & GEN & GEN & -   & \mathrm{THM} & GEN & GEN & GEN & GEN \\
\mathrm{ENR\mbox{-}I} & GEN & GEN & GEN & \mathrm{THM} & -   & GEN & GEN & GEN & GEN \\
\mathrm{CTX\mbox{-}E} & GEN & GEN & GEN & GEN & GEN & -   & \mathrm{THM} & GEN & GEN \\
\mathrm{CTX\mbox{-}S} & GEN & GEN & GEN & GEN & GEN & \mathrm{THM} & -   & GEN & GEN \\
\mathrm{NOR\mbox{-}C} & GEN & GEN & GEN & GEN & GEN & GEN & GEN & -   & \mathrm{THM} \\
\mathrm{NOR\mbox{-}S} & GEN & GEN & GEN & GEN & GEN & GEN & GEN & \mathrm{THM} & - \\
\end{array}
\]

\medskip
For the refined matrices,
$\mathrm{THM}$ indicates separation established by an earlier split theorem.
$\mathrm{GEN\mbox{-}CAR}$ indicates separation established by difference in identity carrier.
All pairs that share an identity carrier are separated by the split theorems,
and no additional separation by identity basis is required for the derived regime set.
Therefore, the matrix contains no $\mathrm{GEN\mbox{-}BAS}$ entries.

\medskip
\noindent\textbf{Refined pairwise non-collapse matrix, Part I.}

\[
\begin{array}{@{}c|ccccc@{}}
 & \mathrm{OBL} & \mathrm{OCC} & \mathrm{REC} & \mathrm{ENR\mbox{-}L} & \mathrm{ENR\mbox{-}I} \\
\hline
\mathrm{OBL} & - & \mathrm{GEN\mbox{-}CAR} & \mathrm{GEN\mbox{-}CAR} & \mathrm{GEN\mbox{-}CAR} & \mathrm{GEN\mbox{-}CAR} \\
\mathrm{OCC} & \mathrm{GEN\mbox{-}CAR} & - & \mathrm{GEN\mbox{-}CAR} & \mathrm{GEN\mbox{-}CAR} & \mathrm{GEN\mbox{-}CAR} \\
\mathrm{REC} & \mathrm{GEN\mbox{-}CAR} & \mathrm{GEN\mbox{-}CAR} & - & \mathrm{GEN\mbox{-}CAR} & \mathrm{GEN\mbox{-}CAR} \\
\mathrm{ENR\mbox{-}L} & \mathrm{GEN\mbox{-}CAR} & \mathrm{GEN\mbox{-}CAR} & \mathrm{GEN\mbox{-}CAR} & - & \mathrm{THM} \\
\mathrm{ENR\mbox{-}I} & \mathrm{GEN\mbox{-}CAR} & \mathrm{GEN\mbox{-}CAR} & \mathrm{GEN\mbox{-}CAR} & \mathrm{THM} & - \\
\mathrm{CTX\mbox{-}E} & \mathrm{GEN\mbox{-}CAR} & \mathrm{GEN\mbox{-}CAR} & \mathrm{GEN\mbox{-}CAR} & \mathrm{GEN\mbox{-}CAR} & \mathrm{GEN\mbox{-}CAR} \\
\mathrm{CTX\mbox{-}S} & \mathrm{GEN\mbox{-}CAR} & \mathrm{GEN\mbox{-}CAR} & \mathrm{GEN\mbox{-}CAR} & \mathrm{GEN\mbox{-}CAR} & \mathrm{GEN\mbox{-}CAR} \\
\mathrm{NOR\mbox{-}C} & \mathrm{GEN\mbox{-}CAR} & \mathrm{GEN\mbox{-}CAR} & \mathrm{GEN\mbox{-}CAR} & \mathrm{GEN\mbox{-}CAR} & \mathrm{GEN\mbox{-}CAR} \\
\mathrm{NOR\mbox{-}S} & \mathrm{GEN\mbox{-}CAR} & \mathrm{GEN\mbox{-}CAR} & \mathrm{GEN\mbox{-}CAR} & \mathrm{GEN\mbox{-}CAR} & \mathrm{GEN\mbox{-}CAR} \\
\end{array}
\]

\medskip
\noindent\textbf{Refined pairwise non-collapse matrix, Part II.}

\[
\begin{array}{@{}c|cccc@{}}
 & \mathrm{CTX\mbox{-}E} & \mathrm{CTX\mbox{-}S} & \mathrm{NOR\mbox{-}C} & \mathrm{NOR\mbox{-}S} \\
\hline
\mathrm{OBL} & \mathrm{GEN\mbox{-}CAR} & \mathrm{GEN\mbox{-}CAR} & \mathrm{GEN\mbox{-}CAR} & \mathrm{GEN\mbox{-}CAR} \\
\mathrm{OCC} & \mathrm{GEN\mbox{-}CAR} & \mathrm{GEN\mbox{-}CAR} & \mathrm{GEN\mbox{-}CAR} & \mathrm{GEN\mbox{-}CAR} \\
\mathrm{REC} & \mathrm{GEN\mbox{-}CAR} & \mathrm{GEN\mbox{-}CAR} & \mathrm{GEN\mbox{-}CAR} & \mathrm{GEN\mbox{-}CAR} \\
\mathrm{ENR\mbox{-}L} & \mathrm{GEN\mbox{-}CAR} & \mathrm{GEN\mbox{-}CAR} & \mathrm{GEN\mbox{-}CAR} & \mathrm{GEN\mbox{-}CAR} \\
\mathrm{ENR\mbox{-}I} & \mathrm{GEN\mbox{-}CAR} & \mathrm{GEN\mbox{-}CAR} & \mathrm{GEN\mbox{-}CAR} & \mathrm{GEN\mbox{-}CAR} \\
\mathrm{CTX\mbox{-}E} & - & \mathrm{THM} & \mathrm{GEN\mbox{-}CAR} & \mathrm{GEN\mbox{-}CAR} \\
\mathrm{CTX\mbox{-}S} & \mathrm{THM} & - & \mathrm{GEN\mbox{-}CAR} & \mathrm{GEN\mbox{-}CAR} \\
\mathrm{NOR\mbox{-}C} & \mathrm{GEN\mbox{-}CAR} & \mathrm{GEN\mbox{-}CAR} & - & \mathrm{THM} \\
\mathrm{NOR\mbox{-}S} & \mathrm{GEN\mbox{-}CAR} & \mathrm{GEN\mbox{-}CAR} & \mathrm{THM} & - \\
\end{array}
\]

%--------------------------------------------
\subsection{Exhaustiveness of Regime Profiles}
\label{subsec:exhaustiveness}

\begin{lemma}[Profile realizability for admissible ontologies]
\label{lem:profile_realizability}
Let $x \in \mathcal{R}$ be an identity carrier belonging to an admissible ontology.
Then there exists a regime profile $P$ in the derived regime set such that
the identity of $x$ is given by the equivalence class $[x]_P$.
\end{lemma}

\begin{proof}
Because $x$ belongs to an admissible ontology, its identity and persistence
conditions are invariant under admissible extension
(Section~\ref{sec:preliminaries}).
By Definition~\ref{def:induced_identity},
identity under a regime profile is determined by the equivalence relation
induced by transformations classified as $\mathrm{PRS}$.

Since the transformation basis $\mathbb{F}$ is fixed
(Section~\ref{subsec:transformation_basis}),
the identity behavior of $x$ under admissible transformations is fully determined
by its classification over $\mathbb{F}$ together with its identity basis.
By the remark following Definition~\ref{def:induced_identity},
no additional structure contributes to identity.

If this transformation-classification pattern is already realized by one of the
previously established regime profiles, then the claim follows immediately.
If not, then by Definition~\ref{def:split_pressure},
failure of realization under a single profile constitutes split pressure.
Such failure must arise from conflicting identity requirements under at least
one transformation family in $\mathbb{F}$.

Sections~\ref{subsec:enr_split}, \ref{subsec:ctx_split}, and \ref{subsec:nor_split}
analyze all transformation families that generate such conflicts for
identity carriers within the structural scope of this analysis,
and establish the corresponding forced refinements.

Therefore, every admissible identity behavior considered in this paper
is realized by some regime profile in the derived regime set,
and the identity of $x$ is given by the corresponding equivalence class $[x]_P$.
\end{proof}

%--------------------------------------------
\subsection{Uniqueness of Regime Realization}
\label{subsec:uniqueness}

\begin{lemma}[Non-overlap of regime realization]
\label{lem:profile_uniqueness}
No admissible identity carrier in the present analysis has identity behavior
realized by more than one regime profile in the derived regime set.
\end{lemma}

\begin{proof}
Suppose $x$ is realized by two regime profiles $P$ and $Q$
in the derived regime set.
Then the identity of $x$ is given by both $[x]_P$ and $[x]_Q$,
so the induced identity relations coincide on the identity behavior of $x$.

If $P$ and $Q$ were distinct realizations of the same identity behavior,
then their induced equivalence relations would coincide,
i.e., $\sim_P = \sim_Q$ on the relevant domain.
By Definition~\ref{def:profile_equivalence}, this implies that $P$ and $Q$ collapse.

However, pairwise non-collapse for all profiles in the derived regime set
is established in Section~\ref{subsec:pairwise_noncollapse}.
Therefore no two profiles in the derived regime set collapse,
and no identity carrier is realized by more than one profile.

Therefore, no admissible identity carrier is realized by more than one regime profile.
\end{proof}

%--------------------------------------------
\subsection{Regime Determinacy}
\label{subsec:determinacy}

\begin{lemma}[Regime determinacy]
\label{lem:regime_determinacy}
For every identity carrier $x$ in an admissible ontology,
there exists exactly one regime profile $P$
such that the identity of $x$ is given by the equivalence class $[x]_P$.
\end{lemma}

\begin{proof}
Existence follows from Lemma~\ref{lem:profile_realizability}.
Uniqueness follows from Lemma~\ref{lem:profile_uniqueness}.
\end{proof}

%--------------------------------------------
\subsection{Representation Theorem}
\label{subsec:representation_theorem}

\begin{theorem}[Neutral representation theorem]
\label{thm:representation}
Let $\mathcal{S}$ be an admissible ontology over the
representation space $\mathcal{R}$
with transformation basis $\mathbb{F}$.
Assume identity carriers in $\mathcal{S}$ are
among the structural forms
considered in Sections~\ref{subsec:enr_split}--\ref{subsec:nor_split}.

Then for every identity carrier $x \in \mathcal{S}$,
there exists exactly one regime profile $P$ in the derived regime set
such that the identity of $x$ is given by the equivalence class $[x]_P$.
\end{theorem}

\begin{proof}
Because $\mathcal{S}$ is admissible,
identity and persistence conditions for each identity carrier
are invariant under admissible extensions
(Section~\ref{sec:preliminaries}).

By Lemma~\ref{lem:profile_realizability},
the identity behavior of each $x \in \mathcal{S}$
is realized by some regime profile $P$
in the derived regime set.

By Lemma~\ref{lem:profile_uniqueness},
no identity carrier is realized by more than one regime profile.

Therefore, each identity carrier in $\mathcal{S}$
admits representation by exactly one regime profile,
with identity given by the equivalence class $[x]_P$.
\end{proof}


The theorem establishes representational adequacy of the derived regime set
for admissible ontologies within the formal setting of this paper.
It does not claim that arbitrary ontological systems admit such a representation.


